% 06-elective-courses.tex -- Advanced Scientific Training
\chapter{Advanced Scientific Training}
\label{ch:elective-courses}

Eighteen elective courses allow doctoral candidates to tailor their training to
the demands of their individual research projects and career goals.  Courses are
grouped into five work-package-aligned clusters.  Each course carries between
three and four ECTS credits.

% ====================================================================
\section{Course Overview}
\label{sec:elective-overview}

\begin{table}[h]
\centering
\small
\begin{tabularx}{\textwidth}{@{}l X l l r l@{}}
  \toprule
  \textbf{Code} & \textbf{Title} & \textbf{WP} & \textbf{Provider} & \textbf{ECTS} & \textbf{Month} \\
  \midrule
  AST-01 & Synthetic Data Generation for Finance                    & WP\,1 & ARC & 4 & M12 \\
  AST-02 & Anomaly Detection in Big Data                            & WP\,1 & BBU & 4 & M18 \\
  AST-03 & Natural Language Processing with Transformers             & WP\,1 & ARC & 4 & M24 \\
  AST-04 & Dependence Structures in High-Frequency Financial Data   & WP\,1 & ASE & 3 & M30 \\
  AST-05 & Reinforcement Learning in Digital Finance                & WP\,2 & UTW & 4 & M12 \\
  AST-06 & Machine Learning in Industry                             & WP\,2 & CAR & 4 & M18 \\
  AST-07 & Deep Learning for Finance                                & WP\,2 & BBU & 3 & M24 \\
  AST-08 & Data-Centric AI                                          & WP\,2 & WWU & 3 & M30 \\
  AST-09 & Cybersecurity in Digital Finance                         & WP\,3 & UTW & 3 & M12 \\
  AST-10 & AI Design in Digital Finance                             & WP\,3 & ASE & 4 & M18 \\
  AST-11 & Barriers in Digital Finance Adoption                     & WP\,3 & WWU & 3 & M24 \\
  AST-12 & Explainable AI in Finance                                & WP\,3 & BFH & 4 & M30 \\
  AST-13 & Digital Finance Regulation                               & WP\,4 & ECB & 3 & M12 \\
  AST-14 & History and Prospects of Digital Finance                  & WP\,4 & UNA & 3 & M18 \\
  AST-15 & Blockchains in Digital Finance                           & WP\,4 & ASE & 4 & M24 \\
  AST-16 & Digital EIT Summer School                                & WP\,5 & EIT & 4 & M18 \\
  AST-17 & Green Digital Finance                                    & WP\,5 & KUT & 3 & M24 \\
  AST-18 & Multi-Criteria Decision Making in Sustainable Finance    & WP\,5 & FRA & 3 & M30 \\
  \bottomrule
\end{tabularx}
\caption{Advanced Scientific Training --- complete course list (from Grant Agreement Table~1.3.c).}
\label{tab:elective-list}
\end{table}

% ====================================================================
\section{Course Descriptions by WP Cluster}
\label{sec:elective-descriptions}

\begin{multicols}{2}

% -----------------------------------------------
\subsection*{WP\,1 --- Data and Financial Data Spaces}

\textbf{AST-01: Synthetic Data Generation for Finance.}\quad
Participants learn to apply deep learning techniques such as generative
adversarial networks to produce synthetic financial data that is
statistically indistinguishable from real data.  Use cases include AI
training for fraud detection and crisis simulation.

\textbf{AST-02: Anomaly Detection in Big Data.}\quad
Students study principles for detecting anomalies in large-scale financial
datasets, including handling data errors through human inspection, outlier
removal, and AI-based gap filling.  The course covers systematic mapping of
data quality issues.

\textbf{AST-03: Natural Language Processing with Transformers.}\quad
Participants combine computational linguistics and role-based modelling of
human language with statistical machine learning and deep learning.  The
course emphasizes deploying advanced transformer architectures for complex
natural language processing tasks in financial contexts.

\textbf{AST-04: Dependence Structures in High-Frequency Financial Data.}\quad
Students learn automatic detection of dependencies between financial vectors,
covering techniques for time-dependent trends, volatility clustering,
seasonality, and fat tails.  Methods include copulas and spectral measures.

% -----------------------------------------------
\subsection*{WP\,2 --- AI for Financial Markets}

\textbf{AST-05: Reinforcement Learning in Digital Finance.}\quad
Participants select and deploy reinforcement learning algorithms relevant to
digital finance, with applications in trading strategy optimization, risk
management, and fraud detection decision-making.

\textbf{AST-06: Machine Learning in Industry.}\quad
Students examine the principles of deploying machine learning in industrial
settings, including business assessment of automation decisions, practical
implementation challenges, and the availability and costs of high-quality data.

\textbf{AST-07: Deep Learning for Finance.}\quad
Participants build and train deep neural networks, identify key architecture
parameters, and implement vectorized neural networks.  The course covers
variance analysis for deep learning applications in financial modelling.

\textbf{AST-08: Data-Centric AI.}\quad
Students learn to empower SMEs in digital finance to deploy AI with limited
datasets by constructing high-quality samples that maximise training impact
and identifying weak spots in data quality.

% -----------------------------------------------
\subsection*{WP\,3 --- Explainable and Fair AI}

\textbf{AST-09: Cybersecurity in Digital Finance.}\quad
Participants study cybersecurity from social behaviour, software, and hardware
perspectives, covering the security of cloud services, EU regulatory compliance,
and techniques for detecting and preventing fraud in financial systems.

\columnbreak

\textbf{AST-10: AI Design in Digital Finance.}\quad
Students survey contemporary AI techniques in digital finance and learn to
design impactful AI systems with consideration for energy consumption, bias,
explainability, and fairness.

\textbf{AST-11: Barriers in Digital Finance Adoption.}\quad
Participants examine hurdles for society-wide adoption of digital finance,
including design principles to include genders, minorities, and vulnerable
groups.  The course addresses the fast-paced start-up climate and competition
dynamics of the industry.

\textbf{AST-12: Explainable AI in Finance.}\quad
Students classify white-box and black-box models, assess the applicability of
classical XAI techniques in finance, and produce audience-dependent
explanations.  The course also covers emerging XAI techniques.

% -----------------------------------------------
\subsection*{WP\,4 --- Blockchain and Digital Innovation}

\textbf{AST-13: Digital Finance Regulation.}\quad
Participants gain an overview of the regulatory field in digital finance,
including an outlook on pending EU regulation changes and best practices for
compliance and monitoring.

\textbf{AST-14: History and Prospects of Digital Finance.}\quad
Students explore past developments in digital finance---digital assets,
algorithmic trading, AI---and consider trends for the next decade, with
reflections on decentralisation and AI-driven transformation.

\textbf{AST-15: Blockchains in Digital Finance.}\quad
Participants study the technical, financial, and legislative principles of
blockchain technology and its applications in digital finance, including the
impact of decentralised finance on traditional financial systems.

% -----------------------------------------------
\subsection*{WP\,5 --- Sustainability of Digital Finance}

\textbf{AST-16: Digital EIT Summer School.}\quad
Students explore how digital technologies disrupt finance and their impact on
society, studying the latest advances through case studies and learning to
write a business plan for fintech ventures.

\textbf{AST-17: Green Digital Finance.}\quad
Participants develop awareness of the energy consumption and ecological footprint
of digital finance, studying techniques for energy-efficient algorithm training
and deployment and the trade-offs between performance and environmental impact.

\textbf{AST-18: Multi-Criteria Decision Making in Sustainable Finance.}\quad
Students learn principles of multi-criteria decision making, including fuzzy
set theory and the analytic hierarchy process, with applications to sustainable
finance decisions.

\end{multicols}
